\section{讨论与局限}\label{sec:discussion}

\subsection{固定宪制与适应性制度}\label{sec:discussion:p2}

本架构刻意不采用无限制自我修改。constitution、transition topology、
capability matrix 与 fixed verifier 不在 runtime 演化。这限制了系统:
真正的新 role 或 transition 需要人类修正。但同一限制也使每项制度变更
可以解释。candidate 能在已知 authority envelope 内改进,却不能重定义
评判自己的 envelope。

宪制本身仍可能错误。hash pinning 使修正显式,却不能保证 pinned rule
明智。因此,人类 review、rationale 与 regression evidence 仍属于 trusted
process。Constitutional ARI-RQGM 减少授予 model output 的权限,但不会
消除定义法律者的责任。

紧急隔离现已保持纪元不变量。T16 强制关闭当前纪元，在隔离违规组件的
同时原子地开启新指纹。这使重放中的遏制点明确，并防止一个指纹代表两个
服务集合。它仍未解决裁决拒绝服务，也不在紧急事务中提供替代角色。

\subsection{权限接缝处的指令遵循}\label{sec:discussion:seams}

最危险的故障多发生在 model-generated record 跨入 deterministic state
之处。schema validation 本身不够:格式完美的记录也可能由错误 role 撰写、
请求 illegal edge、扩张 capability、引用 sealed text,或缺少 repair 所需的
provenance。因此 kernel 分别检查 structure、author、timing、authority、
hash 与 context。

另有一种 semantic limitation:valid attack 或 judgment 仍可能错误。role
separation 与 replay 使错误可观察、可争议,却不会把 model judgment 变成
ground truth。有 execution evidence 或 curated label 时使用 fixed anchor;
其他结论仍是概率性的,不应升级为宪制事实。

同样,“治理失败时研究继续”是有条件的可用性策略,不是一般安全原则。
继续运行应限于沙箱内的候选生成和本地计算。写入外部服务、操作物理设备
或修改生产数据需要独立的固定许可闸门,并在治理故障时关闭。当前内核不
强制这一环境条件。

\subsection{扩展极限}\label{sec:discussion:scaling}

治理增加 model call、replay work、storage 与 boundary latency。budget
可以降低 governance level 或抑制 optional co-evolution,但严重降级的运行
收集到的制度证据也更少。大型 registry 还会增加 candidate、role
interaction 与 repair 必须检查的 dependency edge 数量。

第二版审计摘要绑定类型、标识、事务、载荷和前驱，能揭示改变重放语义的
中间修改，却不能检测历史和本地固定值同时回退到相同旧
前缀。更强的回退抵抗需要外部保存的签名日志头、定期提交给独立见证者,
或单调计数器。永久只追加保存也与机密信息、个人数据和删除要求冲突。
若无销毁密钥式删除或可审计的遮盖记录,不应把敏感明文写入台账。

当前实现假定每个 checkpoint 有一个 authoritative event order。distributed
executor 需要 replicated log、deterministic conflict resolution,以及并发
artifact 产生时究竟哪个 epoch 正在 serving 的明确定义。现有
single-writer discipline 是稳健的 local boundary,还不是 distributed
consensus protocol。

数据同样构成限制。reviewer improvement 需要 held-out anchor;
clean-room regeneration 需要足够抽象的 failure summary;utility-policy
比较需要 stable axis basis。缺少这些输入时,系统保守降级而不虚构优越性
主张。保守降级维护 integrity,却减少适应机会。

反复用同一验证攻击池选择候选和效用策略，可能对已知样例过拟合。可信研究
必须分离开发、选择与秘密最终样例，处理适应性重复选择，并比较同模型、
异模型和异提供者的角色分配。记录级角色分离并不使共享模型系列的错误相互
独立；串谋、裁决拒绝服务和偏置锚点仍是待实证的威胁。

\subsection{展望}\label{sec:discussion:outlook}

下一项实证工作是跨 governance configuration、task family、model 与 seed
的预注册比较。下一项系统工作是显式报告 evidence coverage:对每个 role,
列出实际演练过的 sanction、replay case、anchor 与 successor stage。
重复 policy rewrite 与制度 cycling 需要更长运行才能研究。

优先事项是:外部固定审计日志头，并在进程或操作系统边界隔离角色与检查点
权限；比较固定制度、原
RQGM 和完整系统并做主要组件消融;在带标签的因果图上测量归因与修复精度;
以及量化稿件闸门覆盖率和治理所需时间、token 与费用。

更广的问题是,固定宪制能否支持 open-ended 制度改进,同时避免弱到无法
约束或僵硬到不许改进。Constitutional ARI-RQGM 提供一个具体、可重放的
平台来测量该问题。当前主张刻意更窄:model-defined research role 可以
协同演化,而 registry authority、amendment、provenance 与 final evidence
verification 仍处于其控制之外。
